% 光的干涉（高中）
% keys 波的叠加
% license Usr
% type Tutor

\begin{issues}
\issueTODO 光程和普通机械波的波程关系
\end{issues}

\pentry{机械波及叠加（高中）\nref{nod_mcwave}}{nod_77a3}

在\autoref{sub_mcwave_1}，我们导出了如下规律，倘若两个简谐波满足干涉条件，则在空间中相位差为$2k\pi(k\in \mathbb N_0)$的点是振动加强点，相位差为$(2k+1)\pi$的点则是振动减弱点。接下来，我们将利用这个规律求得双光干涉实验的一些特殊结论。

首先，我们来看一下光的构成。作为电磁波，有电场和磁场两个分量，且都在做简谐振动，振动方向与传播方向垂直，如下图所示。
\begin{figure}[ht]
\centering
\includegraphics[width=12cm]{./figures/add5a31d3e87b27a.png}
\caption{光的电磁场分量} \label{fig_interf_2}
\end{figure}
无论是人眼的光电化学反应，或是其他光电检测设备，所能检测的都是光的强度（即光的能流密度）。可以证明，光强正比于电场振幅平方在一段时间的平均值。又因为电场在做简谐振动，所以要知道双光干涉在某一位置的总光强，我们只需观察电场这一简谐波在此处的合振动振幅平方相对于时间的平均值即可。

由\autoref{eq_mcwave_3} 得，合振幅的平方为
\begin{equation}
C^2=  A^2+B^2+2 A B \cos \left(\phi_1-\phi_2\right)~,
\end{equation}
因此，对于满足干涉条件，即相位差一致的两束光波，合振幅是一常数，大小由相位差决定。又由\autoref{sub_area_2} 知，常数物理量的平均值为该常数。所以空间中波程差为$n\lambda$的位置是振动加强点，波程差为$(2n+1)\lambda/2$则为振动减弱点。\textbf{在光屏上，振动加强的点对应亮条纹，振动减弱的点对应暗条纹。}
\subsection{双缝干涉}
\subsubsection{理论推导}
实验通常利用装有滤光片的单缝，使得一束有各种相位的光源在经过单缝后，成为\textbf{相位和角频率相同的新光源}，由此满足简谐波干涉的条件。

在学习机械波时，我们知道，波实际上传播的是相位，如果知道波的初始振动形式，也就知道其路径上任意一点的振动形式。双缝的位置相对于单缝对称，这样光源到达双缝的时间一致，相位差自然为$0$了。根据我们在“机械波及其叠加”一节中推导出来的结论，对于初始相位差为0的简谐波，亮纹对应的波程差为偶数倍半波长，暗纹对应奇数倍半波长。所以我们只需要用装置参数表示波程差，即可得到亮纹位置，以及其他特点和实验参数的关系。

\begin{figure}[ht]
\centering
\includegraphics[width=8cm]{./figures/dd4445028a4eff52.png}
\caption{双缝干涉原理图} \label{fig_interf_3}
\end{figure}

如上图所示，$x$是条纹所在的位置。考虑到双缝距离$d$远小于屏距$$
\subsubsection{实验装置}
一个经典的双缝干涉实验装置如下图所示。光先通过单缝，成为柱面波，我们可以画出一系列同心圆，每一个圆代表一个相位，最前面的波面是光源的初始相位。也就是说，柱面波到达双缝的必是光源的初始相位。由此，我们在双缝处得到两个相位差为0的光源。
%缺同心圆的图。
\begin{figure}[ht]
\centering
\includegraphics[width=12cm]{./figures/17f92705aa639da5.png}
\caption{双缝干涉实验装置} \label{fig_interf_4}
\end{figure}

\subsection{薄膜干涉}